\documentclass[11pt]{article}
\usepackage[margin=1in]{geometry}
\usepackage{amsmath, amssymb, amsthm}
\usepackage{algorithm}
\usepackage{algpseudocode}
\usepackage{booktabs}
\usepackage{hyperref}

\newtheorem{theorem}{Theorem}
\newtheorem{lemma}[theorem]{Lemma}
\newtheorem{definition}{Definition}
\newtheorem{property}{Property}

\title{Mathematical Formulations, Saturation Dynamics, and Vanishing Gradient Analysis of Sigmoid and Tanh Activations (ALGO-NN-06)}
\author{Algorithms Engineering Group}
\date{\today}

\begin{document}

\maketitle

\begin{abstract}
This document establishes the mathematical specifications, differential calculus, asymptotic saturation properties, and vanishing gradient dynamics of S-shaped bounded activations (Sigmoid, Tanh, Log-Sigmoid, Hard-Sigmoid, and Hard-Tanh). We derive maximum derivative bounds, prove the exponential decay of backpropagated gradients in deep architectures, formulate numerically stable log-domain evaluations, and provide a fully computed numeric trace matching the production implementation contract.
\end{abstract}

\section{Notation}
\label{sec:notation}

Let $\mathbb{R}$ denote the set of real numbers and $\mathbb{N} = \{1, 2, \dots\}$. Matrices are denoted by uppercase bold letters ($\mathbf{X}, \mathbf{Y}, \mathbf{G}$), and scalar values by standard italics.

\begin{itemize}
    \item $N \in \mathbb{N}$: Batch size (number of tensor rows).
    \item $D \in \mathbb{N}$: Feature dimensionality (number of tensor columns).
    \item $\mathbf{X} \in \mathbb{R}^{N \times D}$: Input pre-activation tensor.
    \item $\mathbf{Y} \in \mathbb{R}^{N \times D}$: Activated output tensor.
    \item $\mathbf{G} \in \mathbb{R}^{N \times D}$: Analytic first derivative tensor.
    \item $\sigma(x) = \frac{1}{1 + e^{-x}}$: Standard logistic sigmoid function mapping $\mathbb{R} \to (0, 1)$.
    \item $\tanh(x) = \frac{e^x - e^{-x}}{e^x + e^{-x}} = 2\sigma(2x) - 1$: Hyperbolic tangent mapping $\mathbb{R} \to (-1, 1)$.
    \item $\tau \in \mathbb{R}_{>0}$: Saturation derivative threshold (e.g., $\tau = 0.01$).
\end{itemize}

\section{Problem Definition}
\label{sec:problem_definition}

\begin{definition}[S-Shaped Activation Operators]
For a scalar input $x \in \mathbb{R}$, the component-wise activation functions and their analytic derivatives are defined as follows:

\begin{enumerate}
    \item \textbf{Logistic Sigmoid:}
    \begin{align}
    \sigma(x) &= \frac{1}{1 + e^{-x}}, \label{eq:sigmoid} \\
    \sigma'(x) &= \sigma(x)(1 - \sigma(x)). \label{eq:sigmoid_grad}
    \end{align}

    \item \textbf{Hyperbolic Tangent (Tanh):}
    \begin{align}
    \tanh(x) &= \frac{e^x - e^{-x}}{e^x + e^{-x}}, \label{eq:tanh} \\
    \tanh'(x) &= 1 - \tanh^2(x) = \text{sech}^2(x). \label{eq:tanh_grad}
    \end{align}

    \item \textbf{Log-Sigmoid (Numerically Stable):}
    \begin{align}
    \text{log\_sigmoid}(x) &= \ln(\sigma(x)) = -\ln(1 + e^{-x}) = \begin{cases} -\ln(1 + e^{-x}) & \text{if } x \ge 0, \\ x - \ln(1 + e^x) & \text{if } x < 0, \end{cases} \label{eq:log_sigmoid} \\
    \frac{d}{dx}\text{log\_sigmoid}(x) &= 1 - \sigma(x) = \frac{1}{1 + e^x}. \label{eq:log_sigmoid_grad}
    \end{align}

    \item \textbf{Hard Sigmoid and Hard Tanh (Piecewise Linear):}
    \begin{align}
    \text{hard\_sigmoid}(x) &= \text{clip}(0.2x + 0.5, 0, 1), \quad \text{hard\_sigmoid}'(x) = \begin{cases} 0.2 & \text{if } -2.5 < x < 2.5, \\ 0 & \text{otherwise}. \end{cases} \\
    \text{hard\_tanh}(x) &= \text{clip}(x, -1, 1), \quad \text{hard\_tanh}'(x) = \begin{cases} 1 & \text{if } -1 < x < 1, \\ 0 & \text{otherwise}. \end{cases}
    \end{align}
\end{enumerate}
\end{definition}

\section{Algorithm}
\label{sec:algorithm}

Algorithm~\ref{alg:sigmoid_tanh} specifies the deterministic tensor evaluation, numerical stabilization, and saturation metric calculations.

\begin{algorithm}[H]
\caption{Sigmoid and Tanh Evaluation with Saturation Diagnostics}
\label{alg:sigmoid_tanh}
\begin{algorithmic}[1]
\Require Input tensor $\mathbf{X} \in \mathbb{R}^{N \times D}$, variant $\in \{\text{sigmoid}, \text{tanh}, \text{log\_sigmoid}, \dots\}$, threshold $\tau > 0$.
\Ensure Output tensor $\mathbf{Y} \in \mathbb{R}^{N \times D}$, gradient tensor $\mathbf{G} \in \mathbb{R}^{N \times D}$, saturated fraction $S$, mean gradient $\bar{g}$.
\State $N \gets \text{rows}(\mathbf{X})$, $D \gets \text{cols}(\mathbf{X})$
\State Initialize $\mathbf{Y}, \mathbf{G} \in \mathbb{R}^{N \times D}$
\State $\text{sat\_count} \gets 0$, $\text{grad\_sum} \gets 0.0$
\For{$i = 1 \textbf{ to } N$}
    \For{$j = 1 \textbf{ to } D$}
        \State $x \gets X_{ij}$
        \State $(y, g) \gets \text{EvaluateBounded}(x, \text{variant})$
        \State $Y_{ij} \gets y$, $G_{ij} \gets g$
        \State $\text{grad\_sum} \gets \text{grad\_sum} + g$
        \If{$|g| < \tau$}
            \State $\text{sat\_count} \gets \text{sat\_count} + 1$
        \EndIf
    \EndFor
\EndFor
\State $S \gets \frac{\text{sat\_count}}{N \cdot D}$, $\bar{g} \gets \frac{\text{grad\_sum}}{N \cdot D}$
\State \Return $(\mathbf{Y}, \mathbf{G}, S, \bar{g}, [N, D])$
\end{algorithmic}
\end{algorithm}

\section{Saturation Dynamics and Vanishing Gradients}
\label{sec:correctness}

\begin{theorem}[Maximum Derivative Bounds of S-Curves]
\label{thm:derivative_bounds}
For all $x \in \mathbb{R}$:
\begin{align}
\max_{x \in \mathbb{R}} \sigma'(x) &= \sigma'(0) = \frac{1}{4} = 0.25, \\
\max_{x \in \mathbb{R}} \tanh'(x) &= \tanh'(0) = 1.00.
\end{align}
\end{theorem}

\begin{theorem}[Vanishing Gradient Decay in Deep Sigmoid Networks]
\label{thm:vanishing_gradient}
Consider an $L$-layer feedforward network where each layer computes $\mathbf{h}^{(l)} = \sigma(\mathbf{W}^{(l)} \mathbf{h}^{(l-1)})$. The backpropagated gradient with respect to initial activation $\mathbf{h}^{(0)}$ is given by the Jacobian chain:
\begin{equation}
\frac{\partial \mathcal{L}}{\partial \mathbf{h}^{(0)}} = \frac{\partial \mathcal{L}}{\partial \mathbf{h}^{(L)}} \prod_{l=1}^L \text{diag}(\sigma'(\mathbf{z}^{(l)})) \mathbf{W}^{(l)}.
\end{equation}
Assuming spectral norm $\|\mathbf{W}^{(l)}\|_2 \le 1$, the gradient magnitude satisfies the exponential upper bound:
\begin{equation}
\left\| \frac{\partial \mathcal{L}}{\partial \mathbf{h}^{(0)}} \right\|_2 \le \left\| \frac{\partial \mathcal{L}}{\partial \mathbf{h}^{(L)}} \right\|_2 \prod_{l=1}^L \left\| \text{diag}(\sigma'(\mathbf{z}^{(l)})) \right\|_2 \le \left( \frac{1}{4} \right)^L \left\| \frac{\partial \mathcal{L}}{\partial \mathbf{h}^{(L)}} \right\|_2.
\end{equation}
\end{theorem}

\begin{proof}
By sub-multiplicativity of matrix operator norms:
\begin{equation}
\left\| \frac{\partial \mathcal{L}}{\partial \mathbf{h}^{(0)}} \right\|_2 \le \left\| \frac{\partial \mathcal{L}}{\partial \mathbf{h}^{(L)}} \right\|_2 \prod_{l=1}^L \|\text{diag}(\sigma'(\mathbf{z}^{(l)}))\|_2 \|\mathbf{W}^{(l)}\|_2.
\end{equation}
Since $\max_z \sigma'(z) = 0.25$, each diagonal matrix has spectral norm bounded by $\frac{1}{4}$. With $\|\mathbf{W}^{(l)}\|_2 \le 1$, the product is bounded by $(1/4)^L$. For $L=10$, $(1/4)^{10} \approx 9.5 \times 10^{-7}$, causing total gradient vanishing in early layers.
\end{proof}

\section{Complexity Analysis}
\label{sec:complexity}

\subsection{Time Complexity}
\begin{itemize}
    \item Each element evaluates in $\mathcal{O}(1)$ elementary operations (exponential, division, subtraction).
    \item Total time complexity for tensor $\mathbf{X} \in \mathbb{R}^{N \times D}$ is strictly $\mathcal{O}(N \cdot D)$.
\end{itemize}

\subsection{Space Complexity}
\begin{itemize}
    \item Allocating output activations $\mathbf{Y}$ and local gradients $\mathbf{G}$ requires $2 \cdot N \cdot D$ floating-point entries.
    \item Space complexity is strictly $\mathcal{O}(N \cdot D)$.
\end{itemize}

\section{Worked Example}
\label{sec:worked_example}

Consider input tensor $\mathbf{X} \in \mathbb{R}^{2 \times 2}$ with $\tau = 0.01$:
\begin{equation}
\mathbf{X} = \begin{bmatrix} 0.0 & 2.0 \\ -2.0 & 5.0 \end{bmatrix}.
\end{equation}

\paragraph{1. Logistic Sigmoid ($\sigma(x) = \frac{1}{1 + e^{-x}}$, $\sigma'(x) = \sigma(1 - \sigma)$):}
\begin{align}
x = 0.0: & \quad \sigma(0) = 0.5, \; \sigma'(0) = 0.25, \\
x = 2.0: & \quad \sigma(2) \approx 0.880797, \; \sigma'(2) = 0.880797 \times (1 - 0.880797) \approx 0.104994, \\
x = -2.0: & \quad \sigma(-2) \approx 0.119203, \; \sigma'(-2) = 0.119203 \times (1 - 0.119203) \approx 0.104994, \\
x = 5.0: & \quad \sigma(5) \approx 0.993307, \; \sigma'(5) = 0.993307 \times (1 - 0.993307) \approx 0.006648.
\end{align}

\begin{equation}
\mathbf{Y}_{\text{sigmoid}} = \begin{bmatrix} 0.500000 & 0.880797 \\ 0.119203 & 0.993307 \end{bmatrix}, \quad
\mathbf{G}_{\text{sigmoid}} = \begin{bmatrix} 0.250000 & 0.104994 \\ 0.104994 & 0.006648 \end{bmatrix}.
\end{equation}
At $x = 5.0$, $g = 0.006648 < \tau = 0.01 \implies \text{saturated}$.
$\text{saturated\_fraction} = 1/4 = 0.25$. Mean gradient $\bar{g} = \frac{0.25 + 0.104994 + 0.104994 + 0.006648}{4} \approx 0.116659$.

\section{Numerical Considerations}
\label{sec:numerical}

\begin{enumerate}
    \item \textbf{Logistic Clamping:} To prevent float underflow/overflow in $e^{-x}$, $|x| \ge 40.0$ is mapped to $1.0$ (for $x \ge 40$) and $0.0$ (for $x \le -40$).
    \item \textbf{Stable Log-Sigmoid Formulation:} For negative inputs $x < 0$, $\ln(1 + e^{-x})$ can overflow. Computing $x - \ln(1 + e^x)$ using $\text{log1p}(e^x)$ maintains precision to within machine epsilon ($10^{-16}$).
\end{enumerate}

\begin{thebibliography}{9}

\bibitem{hochreiter1997long}
S.~Hochreiter and J.~Schmidhuber.
\newblock Long short-term memory.
\newblock {\em Neural Computation}, 9(8):1735--1780, 1997.
\newblock DOI: \href{https://doi.org/10.1162/neco.1997.9.8.1735}{10.1162/neco.1997.9.8.1735}.

\bibitem{lecun1998efficient}
Y.~LeCun, L.~Bottou, G.~B.~Orr, and K.-R.~M{\"u}ller.
\newblock Efficient backprop.
\newblock In {\em Neural Networks: Tricks of the Trade}, pages 9--50. Springer, 1998.
\newblock DOI: \href{https://doi.org/10.1007/3-540-49430-8_2}{10.1007/3-540-49430-8\_2}.

\bibitem{glorot2010understanding}
X.~Glorot and Y.~Bengio.
\newblock Understanding the difficulty of training deep feedforward neural networks.
\newblock In {\em Proceedings of the Thirteenth International Conference on Artificial Intelligence and Statistics}, pages 249--256, 2010.

\end{thebibliography}

\end{document}
